\chapter{الدوال والتغيرات}\label{ch:g11:func}

ينظّم هذا الفصل ما ينبغي معرفته عن الدوال قبل
حساب التفاضل: العائلة الصغيرة من \emph{\omterm{def:g11:func:reference}{الدوال المرجعية}} التي يجب
حفظ منحنياتها عن ظهر قلب، والمعنى الدقيق لكون \omterm{def:g11:func:function}{الدالة} \emph{\omterm{def:g11:func:monotone}{متزايدة}} أو
\emph{\omterm{def:g10:functions:variations}{متناقصة}}، وكيف تتصرف التغيرات عند تركيب الدوال.

\section{الدوال والمنحنيات}

\begin{definition}[الدالة ومجموعة التعريف والمنحنى]\label{def:g11:func:function}
\emph{الدالة}\index{دالة} $f$ ترفق بكل عدد $x$ من مجموعة
$D \subseteq \R$ (وهي \emph{مجموعة تعريفها}\index{مجموعة التعريف}) عددًا واحدًا بالضبط
هو $f(x)$. و\emph{منحنيها} (أو تمثيلها البياني) هو مجموعة النقط
$\bigl(x, f(x)\bigr)$ من أجل $x \in D$.
\end{definition}

\begin{example}
تستلزم العبارة $f(x) = \frac{\sqrt{x+2}}{x-1}$ أن يكون $x + 2 \geq 0$ (من أجل
\omterm{def:g11:quad:discriminant}{الجذر} التربيعي) و $x \neq 1$ (لأن المقام غير معدوم): إذن \omterm{def:g11:func:function}{مجموعة التعريف} هي
$\intco{-2}{1} \cup \intoo{1}{+\infty}$.
\end{example}

\section{الدوال المرجعية}

يجب أن تصير الدوال الخمس التالية ومنحنياتها ردود فعل.

\begin{definition}[الدوال المرجعية]\label{def:g11:func:reference}
\index{دوال مرجعية}
\[
x \mapsto x^2 \text{ (مربع)}, \quad
x \mapsto x^3 \text{ (تكعيب)}, \quad
x \mapsto \sqrt{x}, \quad
x \mapsto \frac{1}{x}, \quad
x \mapsto \abs{x}.
\]
ومجموعات تعريفها هي $\R$ و $\R$ و $\intco{0}{+\infty}$ و $\R \setminus \{0\}$
و $\R$ على \omterm{def:g10:coordgeom:system}{الترتيب}. \omterm{def:g10:functions:graph}{ومنحنى} $x \mapsto \frac1x$ هو
\emph{قطع زائد}\index{قطع زائد}.
\end{definition}

\begin{omfigure}
\begin{tikzpicture}
\begin{axis}[omaxis,
  width=7cm, height=5.6cm,
  xmin=-2.3, xmax=2.3, ymin=-3.4, ymax=3.6,
  xtick={-1,1}, ytick={-1,1}]
  \addplot[omDef, thick, domain=-1.85:1.85] {x^2};
  \addplot[omThm, thick, domain=-1.5:1.5] {x^3};
  \node[omDef, font=\small, anchor=east] at (axis cs:-1.5,1.9) {$x^2$};
  \node[omThm, font=\small, anchor=east] at (axis cs:-1.32,-1.73) {$x^3$};
\end{axis}
\end{tikzpicture}%
\begin{tikzpicture}
\begin{axis}[omaxis,
  width=7cm, height=5.6cm,
  xmin=-3.4, xmax=4.4, ymin=-3.4, ymax=3.6,
  xtick={1}, ytick={1}]
  \addplot[omProp, thick, domain=-3.2:-0.31, samples=60] {1/x};
  \addplot[omProp, thick, domain=0.31:4.2, samples=60] {1/x};
  \addplot[omDef, thick, domain=0:4.2] {sqrt(x)};
  \addplot[omThm, thick, domain=-3.2:3.2] {abs(x)};
  \node[omDef, font=\small, anchor=north] at (axis cs:3.6,1.55)
    {$\sqrt{x}$};
  \node[omThm, font=\small, anchor=east] at (axis cs:-2.4,2.8)
    {$\abs{x}$};
  \node[omProp, font=\small, anchor=west] at (axis cs:0.75,2.4)
    {$\tfrac1x$};
\end{axis}
\end{tikzpicture}

{\small المنحنيات المرجعية الخمسة: \omterm{def:g11:quad:parabola}{القطع المكافئ} $x^2$ والتكعيب $x^3$
(إلى اليسار)؛ و $\sqrt x$ و $\abs x$ \omterm{def:g11:func:reference}{والقطع الزائد} $\frac1x$ (إلى اليمين).}
\end{omfigure}

\section{التغيرات على فترة}

\begin{definition}[متزايدة، متناقصة]\label{def:g11:func:monotone}
لتكن $f$ معرَّفة على \omterm{def:g10:numbers:interval}{فترة} $I$. تكون \omterm{def:g11:func:function}{الدالة} $f$
\emph{متزايدة}\index{دالة متزايدة} على $I$ إذا حفظت \omterm{def:g10:coordgeom:system}{الترتيب}:
\[
\text{من أجل كل } u, v \in I: \quad u < v \implies f(u) \leq f(v),
\]
وتكون \emph{\omterm{def:g10:functions:variations}{متناقصة}} إذا قلبت \omterm{def:g10:coordgeom:system}{الترتيب} ($u < v \implies f(u) \geq
f(v)$). ومع المتراجحات التامة في الطرف الأيمن، تكون $f$ متزايدة
\emph{تمامًا} (على \omterm{def:g10:coordgeom:system}{الترتيب} \omterm{def:g10:functions:variations}{متناقصة} تمامًا). \omterm{def:g11:func:function}{والدالة} المتزايدة أو
\omterm{def:g10:functions:variations}{المتناقصة} على $I$ تُسمّى \emph{رتيبة} على $I$.
\end{definition}

\begin{method}[البرهان على تغير انطلاقًا من التعريف]\label{met:g11:func:variation}
خذ عددين كيفيين $u < v$ من \omterm{def:g10:numbers:interval}{الفترة}، وحدّد إشارة
الفرق $f(v) - f(u)$، وذلك عادةً عن طريق \emph{تحليله إلى عوامل}. وإذا كانت
الإشارة ثابتة على \omterm{def:g10:numbers:interval}{الفترة}، فاستنتج.
\end{method}

\begin{proposition}[تغيرات الدوال المرجعية]\label{prop:g11:func:refvariations}
\begin{enumerate}
  \item $x^2$ \omterm{def:g10:functions:variations}{متناقصة} تمامًا على $\intoc{-\infty}{0}$ \omterm{def:g11:func:monotone}{ومتزايدة}
        تمامًا على $\intco{0}{+\infty}$.
  \item $x^3$ \omterm{def:g11:func:monotone}{متزايدة} تمامًا على $\R$.
  \item $\sqrt x$ \omterm{def:g11:func:monotone}{متزايدة} تمامًا على $\intco{0}{+\infty}$.
  \item $\frac1x$ \omterm{def:g10:functions:variations}{متناقصة} تمامًا على $\intoo{-\infty}{0}$ وعلى
        $\intoo{0}{+\infty}$ (لكن \emph{ليس} على اتحادهما).
  \item $\abs x$ \omterm{def:g10:functions:variations}{متناقصة} تمامًا على $\intoc{-\infty}{0}$ و
        \omterm{def:g11:func:monotone}{متزايدة} تمامًا على $\intco{0}{+\infty}$.
\end{enumerate}
\end{proposition}

\begin{proof}
ليكن $u < v$. نحلّل $f(v) - f(u)$ إلى عوامل في كل حالة.

\emph{1.} $v^2 - u^2 = (v-u)(v+u)$. فعلى $\intco{0}{+\infty}$: يكون $v - u > 0$
و $v + u > 0$ (لأن $v > u \geq 0$)، إذن $v^2 > u^2$. وعلى
$\intoc{-\infty}{0}$: يكون $v + u < 0$، إذن $v^2 < u^2$.

\emph{2.} $v^3 - u^3 = (v - u)(v^2 + uv + u^2)$، والعامل الثاني
يساوي $\left(v + \frac u2\right)^2 + \frac{3u^2}{4} > 0$ إلا إذا كان
$u = v = 0$. ومنه $v^3 > u^3$ كلما $v > u$.

\emph{3.} من أجل $0 \leq u < v$:
$\sqrt v - \sqrt u = \frac{(\sqrt v - \sqrt u)(\sqrt v + \sqrt u)}
{\sqrt v + \sqrt u} = \frac{v - u}{\sqrt v + \sqrt u} > 0$.

\emph{4.} من أجل $0 < u < v$:
$\frac1v - \frac1u = \frac{u - v}{uv} < 0$ لأن $u - v < 0$ و $uv > 0$؛
والحساب نفسه من أجل $u < v < 0$. أما على \omterm{def:g10:numbers:interunion}{الاتحاد} فيسقط الادعاء:
إذ $-1 < 1$ لكن $\frac{1}{-1} = -1 < 1 = \frac11$.

\emph{5.} على $\intco{0}{+\infty}$ يكون $\abs x = x$ متزايدًا تمامًا؛ وعلى
$\intoc{-\infty}{0}$ يكون $\abs x = -x$ متناقصًا تمامًا.
\end{proof}

\begin{remark}
النقطة 4 فخّ كلاسيكي: فعبارة ``\omterm{def:g10:functions:variations}{متناقصة} على $\intoo{-\infty}{0}$ وعلى
$\intoo{0}{+\infty}$'' \emph{ليست} كعبارة ``\omterm{def:g10:functions:variations}{متناقصة} على
$\R \setminus \{0\}$''. فالتغيرات لا معنى لها إلا على الفترات.
\end{remark}

\begin{example}[المقارنة دون حساب]
بما أن $\sqrt x$ \omterm{def:g11:func:monotone}{متزايدة} تمامًا، فإن $\sqrt{7} < \sqrt{8}$؛ وبما أن
$x^2$ \omterm{def:g10:functions:variations}{متناقصة} تمامًا على $\intoc{-\infty}{0}$، فإن
$(-3)^2 > (-2)^2$. فالرتابة تحوّل مقارنة الصور إلى
مقارنة السوابق.
\end{example}

\section{دوال جديدة من دوال قديمة}

\begin{proposition}[التحويلات]\label{prop:g11:func:transforms}
لتكن $u$ \omterm{def:g11:func:function}{دالة} معرَّفة على \omterm{def:g10:numbers:interval}{فترة} $I$ وليكن $k \in \R$.
\begin{enumerate}
  \item \omterm{def:g11:func:function}{للدالة} $u + k$ تغيرات $u$ \emph{نفسها} على $I$؛ ومنحنيها
        \omterm{def:g10:functions:graph}{منحنى} $u$ منسحبًا شاقوليًا بمقدار $k$.
  \item ومن أجل $\lambda > 0$، \omterm{def:g11:func:function}{للدالة} $\lambda u$ تغيرات $u$ نفسها؛ ومن أجل
        $\lambda < 0$ تنقلب التغيرات.
  \item وإذا كان $u \geq 0$ على $I$، فللدالة $\sqrt u$ تغيرات
        $u$ نفسها.
  \item وإذا كان $u \neq 0$ وكانت إشارة $u$ ثابتة على $I$، فللدالة $\frac1u$
        تغيرات \emph{معاكسة} لتغيرات $u$.
\end{enumerate}
\end{proposition}

\begin{proof}
ليكن $u < v$ في $I$ (وبتسامح في الترميز، نقارن قيمتَي $x$ حيث $u < v$؛ ونكتب
$s = u(x_1)$ و $t = u(x_2)$ لصورتَي $x_1 < x_2$).

\emph{1.} $(u+k)(x_2) - (u+k)(x_1) = t - s$: فالفرق بين الصورتين
لم يتغير، إذن لم تتغير إشارته.

\emph{2.} $\lambda t - \lambda s = \lambda(t - s)$: فتُحفظ الإشارة إذا كان
$\lambda > 0$، وتنقلب إذا كان $\lambda < 0$.

\emph{3.} إذا كانت $u$ \omterm{def:g11:func:monotone}{متزايدة} و $x_1 < x_2$، فإن $0 \leq s \leq t$، و
بما أن $\sqrt{}$ \omterm{def:g11:func:monotone}{متزايدة} (\cref{prop:g11:func:refvariations}) نجد
$\sqrt s \leq \sqrt t$. والحجة نفسها في حالة التناقص.

\emph{4.} إذا كانت $u$ \omterm{def:g11:func:monotone}{متزايدة}، و $0 < s \leq t$ (أو $s \leq t < 0$)، فإن
$\frac1s \geq \frac1t$ لأن $x \mapsto \frac1x$ \omterm{def:g10:functions:variations}{متناقصة} على كل
جهة من جهتَي $0$: فينقلب \omterm{def:g10:coordgeom:system}{الترتيب}.
\end{proof}

\begin{example}\label{ex:g11:func:composite}
ادرس $f(x) = \sqrt{x^2 + 1}$ على $\R$. \omterm{def:g11:func:function}{الدالة} الداخلية
$u(x) = x^2 + 1$ موجبة، \omterm{def:g10:functions:variations}{ومتناقصة} على $\intoc{-\infty}{0}$ و
\omterm{def:g11:func:monotone}{متزايدة} على $\intco{0}{+\infty}$ (وهي مربع منسحب). وحسب النقطة 3، تكون \omterm{def:g11:func:function}{للدالة} $f$
التغيرات نفسها: \omterm{def:g10:functions:variations}{متناقصة} ثم \omterm{def:g11:func:monotone}{متزايدة}، وقيمتها الصغرى
$f(0) = 1$.
\end{example}

\begin{omfigure}
\begin{tikzpicture}
\begin{axis}[omaxis,
  width=8.6cm, height=5.8cm,
  xmin=-2.6, xmax=2.6, ymin=-1.7, ymax=4.6,
  xtick={-1,1}, ytick={1,2}]
  \addplot[omDef, thick, domain=-1.9:1.9] {x^2};
  \addplot[omThm, thick, domain=-1.55:1.55] {x^2 + 1};
  \addplot[omProp, thick, domain=-2.4:2.4, samples=80]
    {sqrt(x^2 + 1)};
  \node[omDef, font=\small, anchor=west] at (axis cs:1.9,3.4) {$x^2$};
  \node[omThm, font=\small, anchor=west] at (axis cs:1.35,3.3)
    {$x^2{+}1$};
  \node[omProp, font=\small, anchor=north] at (axis cs:-1.9,1.95)
    {$\sqrt{x^2+1}$};
\end{axis}
\end{tikzpicture}

{\small التحويلات في العمل: انسحاب $x^2$ إلى أعلى بمقدار $1$ (بالأحمر) يحفظ
تغيراتها؛ وأخذ \omterm{def:g11:quad:discriminant}{الجذر} التربيعي (بالبرتقالي) يحفظها من جديد، كما في
\cref{ex:g11:func:composite}.}
\end{omfigure}

\section{الدوال الزوجية والفردية}

\begin{definition}[الزوجية]\label{def:g11:func:parity}
لتكن $f$ معرَّفة على \omterm{def:g10:functions:function}{مجموعة تعريف} $D$ متناظرة بالنسبة إلى $0$ (أي أن $-x \in D$
كلما $x \in D$). تكون \omterm{def:g11:func:function}{الدالة} $f$ \emph{زوجية}\index{دالة زوجية}
إذا كان $f(-x) = f(x)$ من أجل كل $x \in D$، وتكون \emph{فردية}\index{دالة فردية} إذا
كان $f(-x) = -f(x)$ من أجل كل $x \in D$.
\end{definition}

\begin{proposition}[المعنى الهندسي]\label{prop:g11:func:paritysym}
\omterm{def:g10:functions:graph}{منحنى} \omterm{def:g11:func:parity}{الدالة الزوجية} متناظر بالنسبة إلى المحور $y$؛ \omterm{def:g10:functions:graph}{ومنحنى}
\omterm{def:g11:func:parity}{الدالة الفردية} متناظر بالنسبة إلى المبدأ.
\end{proposition}

\begin{proof}
إذا كانت $f$ \omterm{def:g11:func:parity}{زوجية} وكانت $(x, y)$ على \omterm{def:g10:functions:graph}{المنحنى} (أي $y = f(x)$)، فإن
$f(-x) = f(x) = y$: فالنقطة المقابلة $(-x, y)$ تقع على \omterm{def:g10:functions:graph}{المنحنى} أيضًا. وإذا كانت
$f$ \omterm{def:g11:func:parity}{فردية}، فإن $f(-x) = -y$: فالنقطة $(-x, -y)$، وهي نظيرة $(x,y)$ بالنسبة إلى
المبدأ، تقع على \omterm{def:g10:functions:graph}{المنحنى}.
\end{proof}

\begin{example}
الدالتان $x^2$ و $\abs x$ زوجيتان؛ والدالتان $x^3$ و $\frac1x$ فرديتان؛
\omterm{def:g11:func:function}{والدالة} $f(x) = x^2 + x$ ليست \omterm{def:g11:func:parity}{زوجية} ولا \omterm{def:g11:func:parity}{فردية}: إذ $f(-1) = 0$ بينما $f(1) = 2$، إذن
$f(-1) \neq \pm f(1)$. \omterm{def:g11:func:parity}{والزوجية} تنصّف العمل: \omterm{def:g11:func:parity}{فالدالة الزوجية} أو \omterm{def:g11:func:parity}{الفردية}
لا تحتاج إلى دراسة إلا على $\intco{0}{+\infty}$.
\end{example}

\section{تمارين}

\begin{exercise}[$\star$]\label{exo:g11:func:1}
أعطِ \omterm{def:g10:functions:function}{مجموعة تعريف} كل \omterm{def:g11:func:function}{دالة}:
\[
f(x) = \sqrt{3 - x}, \qquad
g(x) = \frac{1}{x^2 - 4}, \qquad
h(x) = \frac{\sqrt{x}}{x - 5} .
\]
\end{exercise}

\begin{exercise}[$\star$]\label{exo:g11:func:2}
مستعملًا تغيرات \omterm{def:g11:func:reference}{الدوال المرجعية}، قارن دون آلة
حاسبة:
\[
\sqrt{11} \text{ و } \sqrt{13}; \qquad
(-2.1)^2 \text{ و } (-2.05)^2; \qquad
\frac{1}{0.9} \text{ و } \frac{1}{1.1}; \qquad
(-1.01)^3 \text{ و } (-0.99)^3 .
\]
\end{exercise}

\begin{exercise}[$\star$]\label{exo:g11:func:3}
حدّد هل كل \omterm{def:g11:func:parity}{دالة زوجية} أم \omterm{def:g11:func:parity}{فردية} أم لا هذه ولا تلك:
\[
f(x) = 3x^4 - x^2, \qquad
g(x) = x^3 + x, \qquad
h(x) = x^2 + x^3, \qquad
k(x) = \frac{x}{x^2+1} .
\]
\end{exercise}

\begin{exercise}[$\star$]\label{exo:g11:func:4}
\omterm{def:g11:func:function}{لدالة} $f$ سلوك التغير التالي: \omterm{def:g11:func:monotone}{متزايدة} على
$\intcc{-3}{1}$ من $f(-3) = -2$ إلى $f(1) = 4$، ثم \omterm{def:g10:functions:variations}{متناقصة} على
$\intcc{1}{5}$ نزولًا إلى $f(5) = 0$. كم حلًّا \omterm{def:g10:algebra:equation}{للمعادلة}
$f(x) = 1$ على $\intcc{-3}{5}$؟ وللمعادلة $f(x) = 5$؟
\end{exercise}

\begin{exercise}[$\star\star$]\label{exo:g11:func:5}
مستعملًا \cref{met:g11:func:variation} (إشارة $f(v) - f(u)$)، برهن على أن
$f(x) = x^2 - 4x$ \omterm{def:g11:func:monotone}{متزايدة} تمامًا على $\intco{2}{+\infty}$.
\end{exercise}

\begin{exercise}[$\star\star$]\label{exo:g11:func:6}
دون حساب أي مشتقة، أعطِ تغيرات
\[
f(x) = \sqrt{5 - x} \ \text{ على } \intoc{-\infty}{5},
\quad
g(x) = \frac{1}{x^2 + 1} \ \text{ على } \intco{0}{+\infty},
\quad
h(x) = 3 - 2\sqrt{x} \ \text{ على } \intco{0}{+\infty}.
\]
\end{exercise}

\begin{exercise}[$\star\star$]\label{exo:g11:func:7}
لتكن $f(x) = (x-2)^2 - 3$ على $\R$.
\begin{enumerate}
  \item أعطِ تغيرات $f$ وقيمتها الصغرى.
  \item استنتج تغيرات $g(x) = \frac{1}{(x-2)^2+1}$ وقيمتها
        العظمى. (لاحظ أن $g = \frac{1}{f + 4}$.)
\end{enumerate}
\end{exercise}

\begin{exercise}[$\star\star$]\label{exo:g11:func:8}
يمر \omterm{def:g10:functions:graph}{منحنى} \omterm{def:g11:func:function}{دالة} $f$ بالنقط $(0, 1)$ و $(2, 3)$
و $(4, 1)$، \omterm{def:g11:func:function}{والدالة} $f$ \omterm{def:g11:func:monotone}{متزايدة} على $\intcc{0}{2}$ \omterm{def:g10:functions:variations}{ومتناقصة} على
$\intcc{2}{4}$. ارسم \omterm{def:g10:functions:graph}{منحنى} ممكنًا، ثم ارسم منحنيات
$f + 2$ و $-f$ و $2f$.
\end{exercise}

\begin{exercise}[$\star\star$]\label{exo:g11:func:9}
بيّن أنه يمكن كتابة \omterm{def:g11:func:function}{الدالة} $f(x) = \frac{2x+1}{x+1}$ على الصورة
$f(x) = 2 - \frac{1}{x+1}$، واستنتج تغيراتها على
$\intoo{-1}{+\infty}$.
\end{exercise}

\begin{exercise}[$\star\star$]\label{exo:g11:func:10}
صحيح أم خطأ (برّر أو أعطِ مثالًا مضادًا):
\begin{enumerate}
  \item مجموع دالتين متزايدتين على $I$ متزايد على $I$.
  \item جداء دالتين متزايدتين على $I$ متزايد على
        $I$.
  \item إذا كانت $f$ \omterm{def:g11:func:monotone}{متزايدة} على $I$، فإن $f^2$ ($= f \times f$)
        \omterm{def:g11:func:monotone}{متزايدة} على $I$.
\end{enumerate}
\end{exercise}

\begin{exercise}[$\star\star\star$]\label{exo:g11:func:11}
لتكن $f(x) = x + \frac{1}{x}$ من أجل $x > 0$.
\begin{enumerate}
  \item بيّن أنه من أجل $0 < u < v$،
        $f(v) - f(u) = (v - u)\,\frac{uv - 1}{uv}$.
  \item استنتج أن $f$ \omterm{def:g10:functions:variations}{متناقصة} تمامًا على $\intoc{0}{1}$ و
        \omterm{def:g11:func:monotone}{متزايدة} تمامًا على $\intco{1}{+\infty}$.
  \item استنتج أن $x + \frac1x \geq 2$ من أجل كل $x > 0$، مع تحقق المساواة
        عند $x = 1$ فقط.
\end{enumerate}
\end{exercise}

\section{مسألة: مصنع التناظر}

\begin{problem}[{مسألة نهاية الأسبوع --- كل دالة جزء زوجي
زائد جزء فردي: قواعد الزوجية، ولعبة التحويلات، و
قيم صغرى تُكسب دون حساب تفاضل}]
\label{pb:g11:func:1}
بعض الدوال متناظرة بالنسبة إلى مرآة، وبعضها متناظر
بنصف دورة، ومعظمها لا هذا ولا ذاك --- ومع ذلك تقول معجزة
جبرية صغيرة إن \emph{كل} \omterm{def:g11:func:function}{دالة} هي، بطريقة وحيدة، مجموع
قطعة متناظرة بالمرآة وقطعة متناظرة بنصف دورة. تبرهن هذه
المسألة على المعجزة، وتلعب لعبة التحويلات على المنحنيات
المرجعية (\cref{prop:g11:func:transforms})، وتختم
بتقنية تحليل الفرق إلى عوامل في
\cref{exo:g11:func:11}، التي تكسب مسائل التصغير قبل فصل
كامل من وصول المشتقة.

\medskip
\noindent\textbf{الجزء الأول --- الطلاقة في \omterm{def:g11:func:parity}{الزوجية}.}
\begin{enumerate}
  \item \omterm{def:g11:func:parity}{زوجية} أم \omterm{def:g11:func:parity}{فردية} أم لا هذه ولا تلك (\cref{def:g11:func:parity})؟
        برّر كل حالة: $x^4 - 3x^2$؛ \ $x^3 - x$؛ \ $x^2 + x$؛ \
        $\abs x$؛ \ $\frac1x$.
  \item بيّن المعنى الهندسي لكل حكم
        (\cref{prop:g11:func:paritysym}): أي تناظر \omterm{def:g10:functions:graph}{لمنحنى}
        \omterm{def:g11:func:parity}{دالة زوجية}؟ ولمنحنى \omterm{def:g11:func:parity}{دالة فردية}؟
  \item تتضاعف \omterm{def:g11:func:parity}{الزوجية} كما تتضاعف الإشارات: برهن على أن جداء
        دالتين فرديتين زوجي، وأكمل جدول
        الضرب كله (\omterm{def:g11:func:parity}{زوجية} $\times$ \omterm{def:g11:func:parity}{زوجية}، \omterm{def:g11:func:parity}{وزوجية} $\times$
        \omterm{def:g11:func:parity}{فردية}، \omterm{def:g11:func:parity}{وفردية} $\times$ \omterm{def:g11:func:parity}{فردية}).
  \item نتائج سريعة: إذا كانت $f$ \omterm{def:g11:func:parity}{فردية} ومعرَّفة عند $0$،
        فما $f(0)$؟ وإذا كانت $f$ \omterm{def:g11:func:parity}{زوجية} و $f(3) = 7$، فما
        $f(-3)$؟ وبيّن أن \omterm{def:g11:func:function}{الدالة} الوحيدة \omterm{def:g11:func:parity}{الزوجية} \omterm{def:g11:func:parity}{والفردية} معًا
        هي \omterm{def:g11:func:function}{الدالة} المعدومة.
  \item كم من \omterm{def:g10:functions:graph}{منحنى} \omterm{def:g11:func:parity}{دالة زوجية} يكفي لتحديده
        كله؟ أعطِ \omterm{def:g11:func:function}{دالة} ليست \omterm{def:g11:func:parity}{زوجية} ولا \omterm{def:g11:func:parity}{فردية}، و
        فسّر ما الذي يسقط.
\end{enumerate}

\medskip
\noindent\textbf{الجزء الثاني --- مبرهنة التفكيك.}
من أجل \omterm{def:g11:func:function}{دالة} $f$ معرَّفة على \omterm{def:g10:functions:function}{مجموعة تعريف} متناظرة بالنسبة إلى $0$، ضع
\[
E(x) = \frac{f(x) + f(-x)}{2},
\qquad
O(x) = \frac{f(x) - f(-x)}{2} .
\]
\begin{enumerate}[resume]
  \item احسب $E$ و $O$ من أجل $f(x) = x^3 + x^2 + 1$، و
        تحقق: $E$ \omterm{def:g11:func:parity}{زوجية}، و $O$ \omterm{def:g11:func:parity}{فردية}، و $E + O = f$.
  \item برهن على الصياغة العامة: من أجل أي $f$، تكون \omterm{def:g11:func:function}{الدالة}
        $E$ \omterm{def:g11:func:parity}{زوجية} و $O$ \omterm{def:g11:func:parity}{فردية} و $f = E + O$ --- أي أن \emph{كل
        \omterm{def:g11:func:function}{دالة} تنقسم إلى جزء زوجي وجزء فردي}.
  \item برهن على وحدانية الانقسام: إذا كان
        $f = E_1 + O_1 = E_2 + O_2$ حيث $E_i$ \omterm{def:g11:func:parity}{زوجية} و $O_i$
        \omterm{def:g11:func:parity}{فردية}، فاستعمل السؤال 4 على $E_1 - E_2$.
  \item فكّك $f(x) = \dfrac{1}{1 - x}$ (من أجل
        $x \neq \pm 1$): ضع $E$ و $O$ على المقام
        المشترك $1 - x^2$ وبسّط.
  \item نظرة إلى الأمام، في جملة لكل بند: سيكون الجزآن الزوجي والفردي
        من \omterm{def:g11:func:function}{الدالة} الأسية هما الدالتين التوأمين في السنة الأخيرة
        (جيب التمام الزائدي والجيب الزائدي)؛ وتشغّل
        الغريزة نفسها --- تقسيم كائن إلى قطع
        متناظرة ومعالجة كل قطعة بتناظرها --- تحليلَ
        الصوت والإشارات في الكتب الجامعية.
        فلماذا يكون تقسيم كهذا \emph{مفيدًا} بوجه عام؟
\end{enumerate}

\medskip
\noindent\textbf{الجزء الثالث --- لعبة التحويلات.}
\begin{enumerate}[resume]
  \item انطلاقًا من \omterm{def:g10:functions:graph}{منحنى} $x \mapsto x^2$، صف
        بدقة منحنيات $(x - 3)^2$ و \ $x^2 + 2$ و \
        $(x - 3)^2 + 2$ و \ $-x^2$ و \ $2x^2$
        (\cref{prop:g11:func:transforms}).
  \item ردّ $f(x) = \abs{x - 2} + \abs{x + 2}$ إلى
        عبارة متقطعة على الفترات الثلاث التي يقطعها $-2$
        و $2$، وصف منحنيها (وله قاع مسطح
        شهير).
  \item أي التحويلات تحترم \omterm{def:g11:func:parity}{الزوجية}؟ قرّر، مع
        برهان أو مثال مضاد: الانسحاب الأفقي
        $f(x - 3)$؛ والانسحاب الشاقولي $f(x) + 2$؛ والتمديد الشاقولي
        $2 f(x)$ (وكلها مطبقة على \omterm{def:g11:func:parity}{دالة زوجية} $f$).
  \item حل $\abs{x - 2} + \abs{x + 2} = 6$ مستعملًا
        العبارة المتقطعة في السؤال 12.
  \item الهندسة العكسية: \omterm{def:g11:quad:parabola}{لقطع مكافئ} رأسه $(2, -3)$
        ويمر بالنقطة $(0, 1)$. أعد بناء معادلته على
        \omterm{thm:g11:quad:canonical}{الصورة النموذجية}.
\end{enumerate}

\medskip
\noindent\textbf{الجزء الرابع --- قيم صغرى دون حساب تفاضل.}
\begin{enumerate}[resume]
  \item قابِل \cref{exo:g11:func:11} من أجل
        $g(x) = x + \dfrac4x$ على $\intoo{0}{+\infty}$: بيّن أن
        $g(v) - g(u) = (v - u)\,\dfrac{uv - 4}{uv}$، وحدّد
        \omterm{def:g10:functions:extrema}{القيمة الصغرى}، وأعطِ قيمتها.
  \item طبّق: من بين كل المستطيلات ذات المساحة $16$، أيها له
        أصغر محيط؟ (اكتب $P = 2\left(x +
        \frac{16}{x}\right)$ واستعمل آلة السؤال 16.) ويختم
        صديق قديم.
  \item حدّد تغيرات
        $h(x) = \sqrt{x^2 + 1}$ على $\R$ كلها، جامعًا بين حجة
        تركيب على $\intco{0}{+\infty}$ وحجة
        \omterm{def:g11:func:parity}{زوجية} من أجل الباقي.
  \item أقرب نقطة: أوجد نقطة \omterm{def:g10:functions:graph}{المنحنى}
        $y = \sqrt x$ الأقرب إلى $(3, 0)$. (صغّر
        \emph{مربع} المسافة --- وهو مقدار تربيعي في $x$ --- و
        فسّر لماذا يكون تصغير المربع مشروعًا.)
  \item الخاتمة، وهي جولة في المصنع: المنحنيات المرجعية هي المادة
        الخام؛ والتحويلات هي الآلات؛ \omterm{def:g11:func:parity}{والزوجية} والتفكيك هما
        ضبط الجودة؛ وتحليل الفرق إلى عوامل هو
        منضدة القياس. جملة واحدة عن كل بند --- وسمِّ
        الأداة القوية التي يقدمها الفصل التالي.
\end{enumerate}
\end{problem}
